%% ma: L10-B14-0009
%% lop: 10
%% chuong: 5
%% bai: 14
%% dang: 10.14.2
%% loai: TL
%% muc: TH
%% dap_an: "-"
%% nguon: "chuyen-de-cac-so-dac-trung-cua-mau-so-lieu-khong-ghep-nhom-toan-10-chuong-trinh-moi.pdf (GV. Phan Nhật Linh) (P04), Bài 03 Dạng 1 Bài tập tự luận 9"
%% kien_thuc: "Lớp 10 Bài 13 (tứ phân vị), Bài 14 (khoảng tứ phân vị)"
%% trang_thai: chinh-thuc
%% ghi_chu: "Tự duyệt theo yêu cầu giáo viên 22/10/2026."
\begin{ex}
Sản lượng lúa (đơn vị: tạ) của $40$ thửa ruộng thí nghiệm có cùng diện tích được cho trong bảng sau:
\begin{center}
\begin{tabular}{|c|c|c|c|c|c|}\hline
Sản lượng & $20$ & $21$ & $22$ & $23$ & $24$\\\hline
Tần số & $5$ & $8$ & $11$ & $10$ & $6$\\\hline
\end{tabular}
\end{center}
Tính khoảng biến thiên và khoảng tứ phân vị của mẫu số liệu.
\loigiai{$R=24-20=4$.

Tần số tích lũy $5$; $13$; $24$; $34$; $40$. $Q_1=\dfrac{x_{10}+x_{11}}2=21$; $Q_2=\dfrac{x_{20}+x_{21}}2=22$; $Q_3=\dfrac{x_{30}+x_{31}}2=23$.

$\Delta_Q=23-21=2$.}
\end{ex}
